\documentclass[12pt, a4paper]{article}

% --- Lingua e Codifica ---
\usepackage[utf8]{inputenc}
\usepackage[T1]{fontenc}
\usepackage[italian]{babel}
\usepackage[expansion=false]{microtype}

% --- Impaginazione e Pacchetti ---
\usepackage[margin=2.5cm]{geometry}
\usepackage{amsmath, amssymb, amsfonts, amsthm}
\usepackage{graphicx, hyperref, booktabs}
\usepackage[document]{ragged2e}
\usepackage{xcolor} 

% Titolo e intestazione
\title{Flusso Progetto}
\author{Cavina Mattia}


\begin{document}

Il progetto è iniziato con una riunione preliminare con il Project Manager Claudio Rossi per definire gli obiettivi e le risorse disponibili il 4 dicembre 2025.
Durante questo incontro sono state delineate le fasi principali, le tempistiche e i ruoli chiave coinvolti.

\subsection{Fasi del progetto}
\begin{itemize}
      \item \textbf{Raccolta dei requisiti:} Acquisizione sistematica dei requisiti mediante interviste condotte con gli utenti finali e stakeholder coinvolti collateralmente.
      \item \textbf{Elaborazione e analisi interviste:} Trascrizione e analisi delle interviste per identificare i requisiti funzionali e non funzionali.
      \item \textbf{Analisi requisiti:} Elaborazione e documentazione sistematica dei requisiti raccolti.
      \item \textbf{Progettazione del database:} Conversione, normalizzazione e progettazione relazionale del database Access esistente.
      \item \textbf{Sviluppo software:} Implementazione del software in due fasi: versione iniziale con funzionalità prioritarie, successivamente funzionalità avanzate.
      \item \textbf{Testing:} Test parallelo allo sviluppo, includendo test unitari e test di integrazione del sistema.
      \item \textbf{Test con utenti:} Validazione della prima versione con un gruppo selezionato di utenti.
      \item \textbf{Affiancamento:} Periodo di esercizio parallelo tra il nuovo e il vecchio sistema per consentire l'adattamento graduale degli utenti.
      \item \textbf{Rilascio definitivo:} Passaggio completo al nuovo sistema dopo il periodo di affiancamento.
\end{itemize}

\end{document}